\documentclass{article}

\usepackage{bm}
\usepackage{ctex}
\usepackage{tikz}
\usepackage{float}
\usepackage{xcolor}
\usepackage{booktabs}
\usepackage{setspace}
\usepackage{datetime}
\usepackage{geometry}
\usepackage{graphicx}
\usepackage{enumitem}
\usepackage{tcolorbox}
\usepackage{nicematrix}
\usepackage{amsmath, amssymb, amsfonts}
\usepackage[fontsize = 12pt]{fontsize}

\renewcommand{\thesubsection}{(\alph{subsection})}
\newcommand{\diff}[2]{\dfrac{\mathrm{d}#1}{\mathrm{d}#2}}
\newcommand{\diffn}[3]{\dfrac{\mathrm{d}^{#3}#1}{\mathrm{d}#2^{#3}}}
\newcommand{\piff}[2]{\dfrac{\partial{#1}}{\partial{#2}}}
\newcommand{\eval}[2]{\left.#1\right|_{#2}}
\newcommand{\e}{\mathrm{e}}
\renewcommand{\d}{\mathrm{d}}
\newcommand{\barline}[1]{\overline{#1\rule{0pt}{1.5ex}}}
\newcommand{\mathsize}{\fontsize{6.5pt}{10pt}\selectfont}

\setcounter{MaxMatrixCols}{20}

\geometry{
    a4paper,
    left = 2cm,
    right = 2cm,
    top = 2cm,
    bottom = 2cm
}

\setlist[1]{
    itemsep = 0pt,
    partopsep = 0pt,
    parsep = 0pt,
    topsep = 5pt,
    leftmargin = 2em,
    labelsep = 0.25em
}

\title{图像处理：第六次作业}
\author{Illusionna \quad 2025........004 \quad 人工智能学院}
\date{\currenttime,\ \today}


\begin{document}

\maketitle
\begin{spacing}{1.75}
    \tableofcontents
\end{spacing}
\thispagestyle{empty}
\clearpage
\setcounter{page}{1}



\section{灰度值及其在图像中出现的概率}
\subsection{哈夫曼编码构建最优二叉树}
\begin{table}[H]
    \centering
    \setlength{\tabcolsep}{64pt}
    \renewcommand{\arraystretch}{1.2}
    \caption{一组灰度值及其在图像中出现的频率}
    \begin{tabular}{cc}
        \bottomrule[1.5pt]
        灰度值 & 频率 \\
        \midrule[0.75pt]
        0$^{'}$ & 45 \\
        1$^{'}$ & 35 \\
        2$^{'}$ & 20 \\
        3$^{'}$ & 15 \\
        4$^{'}$ & 15 \\
        5$^{'}$ & 10 \\
        6$^{'}$ & 5 \\
        7$^{'}$ & 5 \\
        \toprule[1.5pt]
    \end{tabular}
\end{table}

\paragraph*{Step 1：按照频率从小到大排序}~{}

\begin{table}[H]
    \centering
    \setlength{\tabcolsep}{64pt}
    \renewcommand{\arraystretch}{1.2}
    \caption{一组灰度值及其在图像中出现的频率}
    \begin{tabular}{cc}
        \bottomrule[1.5pt]
        灰度值 & 频率 \\
        \midrule[0.75pt]
        6$^{'}$ & 5 \\
        7$^{'}$ & 5 \\
        5$^{'}$ & 10 \\
        3$^{'}$ & 15 \\
        4$^{'}$ & 15 \\
        2$^{'}$ & 20 \\
        1$^{'}$ & 35 \\
        0$^{'}$ & 45 \\
        \toprule[1.5pt]
    \end{tabular}
\end{table}

\paragraph*{Step 2：依次取最小频率的两个灰度节点，其权重为两者之和，构建哈夫曼树}~{}

\begin{itemize}
    \item 6$^{'}$ + 7$^{'}$ = N1 = 5 + 5 = 10
    \item 5$^{'}$ + N1 = N2 = 10 + 10 = 20
    \item 3$^{'}$ + 4$^{'}$ = N3 = 15 + 15 = 30
    \item 2$^{'}$ + N2 = N4 = 20 + 20 = 40
    \item N3 + 1$^{'}$ = N5 = 30 + 35 = 65
    \item N4 + 0$^{'}$ = N6 = 40 + 45 = 85
    \item N5 + N6 = N7 = 65 + 85 = 150
\end{itemize}

\paragraph*{Step 3：约定左分支子树编码为 0，右方为 1，构建哈夫曼树}~{}

\begin{figure}[H]
    \centering
    \includegraphics*[scale=0.23]{./figs/huffman.pdf}
    \caption{最优哈夫曼二叉树}
\end{figure}

\paragraph*{Step 4：构建哈夫曼编码表}~{}

\begin{table}[H]
    \centering
    \setlength{\tabcolsep}{32pt}
    \renewcommand{\arraystretch}{1.2}
    \caption{一组灰度值及其在图像中出现的频率}
    \begin{tabular}{cccc}
        \bottomrule[1.5pt]
        灰度值 & 频率 & 哈夫曼编码 & 码长 \\
        \midrule[0.75pt]
        0$^{'}$ & 45 & 11 & 2 \\
        1$^{'}$ & 35 & 01 & 2 \\
        2$^{'}$ & 20 & 100 & 3 \\
        3$^{'}$ & 15 & 000 & 3 \\
        4$^{'}$ & 15 & 001 & 3 \\
        5$^{'}$ & 10 & 1010 & 4 \\
        6$^{'}$ & 5 & 10110 & 5 \\
        7$^{'}$ & 5 & 10111 & 5 \\
        \toprule[1.5pt]
    \end{tabular}
\end{table}

\paragraph*{Step 5：计算平均码长}~{}
\[ \barline{L}=\dfrac{\displaystyle\sum_{i=1}^{n}f_i\times L_i}{\displaystyle\sum_{i=1}^{n}f_i}=\dfrac{45\times2+35\times2+20\times3+\cdots+5\times5+5\times5}{45+35+20+15+15+10+5+5}=\dfrac{400}{150}\approx2.667 \]



\subsection{哈夫曼编码节省的存储空间}
因为总像素（即总频率）为 150，由于每个像素 3 个位数，所以一共 $150\times3=450$，而哈夫曼编码压缩后总大小为 400，所以节省了 50 比特，因此节省比例是 $50\ /\ 450\approx0.111$，即使用哈夫曼编码可以节省 11.11\% 的存储空间。



\subsection{哈夫曼编码在图像压缩中的优劣性}
哈夫曼编码在数字图像压缩中通常扮演最后一步打包工作，而不是唯一的压缩手段，具体的优劣势如下：

\begin{itemize}
    \item 优势 1：无损压缩，哈夫曼编码完全基于统计概率，解码后的数据与编码前的完全一致，没有丢失任何信息；
    \item 优势 2：实现简单且高效，算法逻辑清晰，使用查表法，计算复杂度很低，编码解码速度非常快，适合实时处理；
    \item 优势 3：符号级最优，在不考虑符号间相关性且必须使用整数位来表示一个符号的前提下，哈夫曼编码被证明是效率最高的编码方式；
    \item[-]
    \item 劣势 1：忽略了像素间的相关性，哈夫曼编码通常把像素当作独立的事件来看，但实际情况下，一个蓝色像素点附近的像素也大概率是蓝色，如果直接对原始图像做哈夫曼编码，效率并不高，没有利用这种邻居相似的规律，现代标准，如 JPEG 通常先通过 DCT 变化去除相关性，剩下不相关度数据系数再交给哈夫曼编码；
    \item 劣势 2：整数位限制，哈夫曼编码强制每个符号码长必须是整数，如果符号出现的概率极高，理论上不需要很多位就能表示，但哈夫曼最少也得分配 1 个位，导致浪费，现代通常采用算术编码，逼近分数位，因此比哈夫曼压缩率更高，但计算也更复杂；
    \item 劣势 3：误码敏感性，由于是变长编码且没有固定的分隔符，一旦数据流中有一个比特发生翻转，比如 0 变为 1，解码器可能会“对齐错误”，导致后面一连串数据全部计算错误。
\end{itemize}



\section{算术解码}
对编码模型信息 0.23355 进行解码：

\begin{table}[H]
    \centering
    \setlength{\tabcolsep}{64pt}
    \renewcommand{\arraystretch}{1.2}
    \caption{算术编码与解码}
    \begin{tabular}{cc}
        \bottomrule[1.5pt]
        符号 & 频率 \\
        \midrule[0.75pt]
        a & 0.2 \\
        e & 0.3 \\
        i & 0.1 \\
        o & 0.2 \\
        u & 0.1 \\
        ! & 0.1 \\
        \toprule[1.5pt]
    \end{tabular}
\end{table}

\paragraph*{Step 1：构建累积概率区间}~{}

\begin{itemize}
    \item a：0.2 $\to$ [0, 0.2)
    \item e：0.3 $\to$ [0.2, 0.5)
    \item i：0.1 $\to$ [0.5, 0.6)
    \item o：0.2 $\to$ [0.6, 0.8)
    \item u：0.1 $\to$ [0.8, 0.9)
    \item !：0.1 $\to$ [0.9, 1)
\end{itemize}

\begin{table}[H]
    \centering
    \setlength{\tabcolsep}{32pt}
    \renewcommand{\arraystretch}{1.2}
    \caption{算术编码与解码累计区间}
    \begin{tabular}{ccc}
        \bottomrule[1.5pt]
        符号 & 频率 & 累计区间 \\
        \midrule[0.75pt]
        a & 0.2 & [0, 0.2) \\
        e & 0.3 & [0.2, 0.5) \\
        i & 0.1 & [0.5, 0.6) \\
        o & 0.2 & [0.6, 0.8) \\
        u & 0.1 & [0.8, 0.9) \\
        ! & 0.1 & [0.9, 1) \\
        \toprule[1.5pt]
    \end{tabular}
\end{table}

\paragraph*{Step 2：解码迭代过程计算}~{}

\[ V^{'}=\dfrac{V-\min}{\max-\min} \]

\begin{itemize}[
    itemsep = 10pt,
    partopsep = 0pt,
    parsep = 0pt,
    topsep = 5pt,
    leftmargin = 2em,
    labelsep = 0.25em
]
    \item 输入：0.23355，输出：e，迭代：$\dfrac{0.23355-0.2}{0.5-0.2}\approx0.11183$
    \item 输入：0.11183，输出：a，迭代：$\dfrac{0.11183-0}{0.2-0}\approx0.55915$
    \item 输入：0.55915，输出：i，迭代：$\dfrac{0.55915-0.5}{0.6-0.5}\approx0.59167$
    \item 输入：0.59167，输出：i，迭代：$\dfrac{0.59167-0.5}{0.6-0.5}\approx0.91667$
    \item 输入：0.91667，输出：! 为终止程序，退出，返回结果 eaii!
\end{itemize}



\section{名词解释}
\subsection{信源编码与解码器}
\textbf{定义：}信源编码的主要目的是为了提高通信的有效性。它通过压缩数据、去除信源中的冗余信息，将信源符号转换为适合传输或存储的数字信号，通常是二进制流。常见的信源编码算法包括哈夫曼编码和算术编码。解码器则是编码的逆过程，试图无失真或在允许失真范围内恢复原始信源信息。



\subsection{信道编码与解码器}
\textbf{定义：}信道编码的主要目的是为了提高通信的可靠性。针对信道中可能存在的噪声和干扰，信道编码器在信息序列中有规律地增加一些冗余码元，如校验位，以便在接收端进行检错或纠错。解码器利用这些冗余信息来检测或纠正传输过程中产生的错误，从受损的信号中恢复出原始发送的数据。



\subsection{信息量 / 自信息}
\textbf{定义：}自信息是衡量一个特定事件发生时所带来的信息量，或者是该事件发生前的不确定性。一个事件发生的概率越小，其发生时所包含的信息量就越大。

\textbf{数学公式：}离散随机变量 $X$ 的某个取值 $x$ 发生的概率为 $p(x)$，则其自信息 $I(x)$ 定义：
\[ I(x) = -\log_2 p(x) \]
单位通常为比特。



\subsection{熵}
\textbf{定义：}熵，也称为香农熵，表示信源的平均不确定性，或者是信源输出的平均信息量。它是所有可能事件自信息的数学期望。

\textbf{数学公式：}对于离散随机变量 $X$，其概率分布为 $p(x_i)$，则熵 $H(X)$ 为：
\[ H(X) = \sum_{i} p(x_i) I(x_i) = -\sum_{i} p(x_i) \log_2 p(x_i) \]

\subsection{条件熵}
\textbf{定义：}条件熵 $H(Y|X)$ 表示在已知随机变量 $X$ 的条件下，随机变量 $Y$ 的不确定性。它衡量了在接收到 $X$ 后，关于 $Y$ 仍然存在的平均不确定度，通常用于描述信道中的噪声损失。

\textbf{数学公式：}
\[ H(Y|X) = -\sum_{x} \sum_{y} p(x,y) \log_2 p(y|x) \]



\subsection{互信息}
\textbf{定义：}互信息 $I(X;Y)$ 表示两个随机变量之间的相关性度量。在通信中，它表示接收端在收到信号 $Y$ 后，获得的关于发送端 $X$ 的信息量。也就是“发送 $X$ 前的不确定性”减去“收到 $Y$ 后 $X$ 的不确定性”。

\textbf{数学公式：}
\[ I(X;Y) = H(X) - H(X|Y) = H(Y) - H(Y|X) \]



\subsection{信道容量}
\textbf{定义：}信道容量是指信道能够无差错传输信息的最大理论速率。根据香农第二定理，信道编码定理，只要信息传输速率低于信道容量，就存在一种编码方式，使得误差概率可以任意小。

\textbf{数学公式：}信道容量 $C$ 是互信息 $I(X;Y)$ 关于输入概率分布 $p(x)$ 的最大值：
\[ C = \max_{p(x)} I(X;Y) \]



\section{海明码的编码解码计算方法}
海明码（7，4）是一种能够纠正 1 位错误的线性分组码，它通过在 4 位数据位中加入 3 位校验位，构成总长为 7 位的码字。

\begin{itemize}
    \item 码字总长度：$n = 7$
    \item 数据位长度：$k = 4$
    \item 校验位长度：$m = n - k = 3$
\end{itemize}



\subsection{编码方法}
\paragraph*{位置分配}~{}

设原始数据为 $\mathbf{d} = (d_1, d_2, d_3, d_4)$。在海明码中，校验位 $p_i$ 被放置在 $2$ 的幂次方位置（即第 1, 2, 4 位），数据位按顺序填充剩余位置。生成的码字 $\mathbf{c}$ 结构如下：
\[ \mathbf{c} = (p_1, p_2, d_1, p_3, d_2, d_3, d_4) \]

\paragraph*{校验位计算}~{}

校验位的计算基于模 2 加法，即异或运算，记为 $\oplus$，每个校验位负责检查其对应二进制索引中该位为 1 的所有位置：
\begin{align*}
    p_1 &= d_1 \oplus d_2 \oplus d_4 \quad (\text{校验第 } 1, 3, 5, 7 \text{ 位}) \\
    p_2 &= d_1 \oplus d_3 \oplus d_4 \quad (\text{校验第 } 2, 3, 6, 7 \text{ 位}) \\
    p_3 &= d_2 \oplus d_3 \oplus d_4 \quad (\text{校验第 } 4, 5, 6, 7 \text{ 位})
\end{align*}



\subsection{解码与纠错}
假设接收端收到的码字为 $\mathbf{r} = (r_1, r_2, r_3, r_4, r_5, r_6, r_7)$，为了检测错误，我们需要计算伴随式 $\mathbf{S} = (s_3, s_2, s_1)$。

\paragraph*{伴随式计算}~{}
\begin{align*}
    s_1 &= r_1 \oplus r_3 \oplus r_5 \oplus r_7 \\
    s_2 &= r_2 \oplus r_3 \oplus r_6 \oplus r_7 \\
    s_3 &= r_4 \oplus r_5 \oplus r_6 \oplus r_7
\end{align*}

\paragraph*{错误定位逻辑}~{}

将伴随式向量 $\mathbf{S}$ 视为一个二进制数 $(s_3 s_2 s_1)_{[2]}$：

\begin{itemize}
    \item 若 $\mathbf{S} = (0,0,0)$，表示传输无误。
    \item 若 $\mathbf{S} \neq 0$，将该二进制数转换为十进制整数 $k$，则说明第 $k$ 位发生了错误。
    \item 纠错操作：将接收到的第 $k$ 位数据取反（$0 \to 1$ 或 $1 \to 0$）即可恢复正确数据。
\end{itemize}



\subsection{矩阵表示法}
在计算机实现中，通常使用生成矩阵 $G$ 和校验矩阵 $H$ 进行运算。海明编码过程可以表示为 $\mathbf{c}=\mathbf{d}G$：
\[ G=\left[\begin{matrix}
    1 & 1 & 1 & 0 & 0 & 0 & 0 \\
    1 & 0 & 0 & 1 & 1 & 0 & 0 \\
    0 & 1 & 0 & 1 & 0 & 1 & 0 \\
    1 & 1 & 0 & 1 & 0 & 0 & 1 \\
\end{matrix}\right] \]

使用校验矩阵 $H$ 进行解码：
\[ H=\begin{pmatrix}
    0 & 0 & 0 & 1 & 1 & 1 & 1 \\
    0 & 1 & 1 & 0 & 0 & 1 & 1 \\
    1 & 0 & 1 & 0 & 1 & 0 & 1
\end{pmatrix} \]

伴随式的计算可通过矩阵乘法获得：
\[ \mathbf{S}^{\top} = H \cdot \mathbf{r}^{\top} \pmod 2 \]


\section{证明二元表达式}
证明：
\[ (A\bullet B)^{c}=A^{c}\circ\widehat{B} \]

展开闭运算的定义，将 $A\bullet B$ 替换为其定义公式 LHS，再应用腐蚀的对偶性，根据 $(X\circleddash Y)^{c}=X^{c}\oplus\widehat{Y}$ 对整体求补集，然后应用 $(X\oplus Y)^{c}=X^{c}\circleddash \widehat{Y}$ 膨胀的对偶性，最后利用 $X\circ Y=(X\circleddash Y)\oplus Y$ 开运算的定义，推导出结果 RHS：
\begin{align*}
    (A\bullet B)^{c}&=\left[(A\oplus B)\circleddash B\right]^{c}=\text{LHS}
    \\
    &=(A\oplus B)^{c}\oplus\widehat{B}
    \\
    &=(A^{c}\circleddash\widehat{B})\oplus\widehat{B}
    \\
    &=A^{c}\circ\widehat{B}
\end{align*}



\end{document}